\clearpage
\section{Partially paired ATAC--RNA factorization}
\label{sec:multiview}
Construct the union of cell identifiers. Each modality has its own features,
rank, noise and observation mask:
\[
Y^A=L^A(F^A)^\top+E^A,\qquad
Y^R=L^R(F^R)^\top+E^R.
\]
The joint loading prior is
\begin{align}
p(L_i^A,L_i^R\mid X_i)
={}&\prod_{k=1}^{K_A}g^A_k(L^A_{ik}\mid X_i,L^A_{i,<k})\notag\\
&\times\prod_{k=1}^{K_R}
g^R_k(L^R_{ik}\mid X_i,L_i^A,L^R_{i,<k}).
\label{eq:multiviewprior}
\end{align}
This is the same autoregressive construction with all ATAC nodes ordered
before RNA nodes. It permits different ranks and does not require a common
loading matrix. An ash prior is learned separately for each feature column.
The directed prior describes a conditional association, not an identified
causal direction.

\subsection{Observed likelihood and the three cell groups}
Let \(m_i^A,m_i^R\in\{0,1\}\) indicate whether a modality has any observed
entries in cell \(i\). Writing each modality likelihood using only its
observed features, the row likelihood at fixed feature factors is
\begin{equation}
\int p(L_i^A\mid X_i)p(L_i^R\mid L_i^A,X_i)
\,p(Y_i^{A,\mathrm{obs}}\mid L_i^A)^{m_i^A}
\,p(Y_i^{R,\mathrm{obs}}\mid L_i^R)^{m_i^R}
\,dL_i^A\,dL_i^R.
\label{eq:observedmultiview}
\end{equation}
This conditions on the missingness pattern and assumes it is ignorable for
the model parameters. A mechanism depending on unobserved biology may need
an additional missingness model.

\begin{description}[leftmargin=0pt]
\item[Paired cell.] Both likelihoods contribute. An ATAC update includes
RNA-child feedback; an RNA update averages over uncertain ATAC parents.
\item[RNA-only cell.] Direct ATAC likelihood statistics satisfy
\(a^A=b^A=0\). ATAC loadings remain latent parents of the observed RNA.
Their updates use the prior and child terms in \eqref{eq:qstar}; there is
no division by a zero precision and no substitution of a zero loading.
\item[ATAC-only cell.] The entire unobserved RNA branch integrates out:
\[
\int p_\theta(L_i^R\mid L_i^A,X_i)\,dL_i^R=1.
\]
The cell informs the ATAC model but provides no direct evidence for the
separately parameterized RNA conditional. RNA predictions are generated
after fitting by averaging that conditional over inferred ATAC uncertainty.
\end{description}

\clearpage
\subsection{Why an unobserved child should be collapsed}
Retaining a factorized \(q(L_i^A)q(L_i^R)\) for an ATAC-only cell adds
\begin{equation}
\E_q\log p_\theta(L_i^R\mid L_i^A,X_i)+H(q(L_i^R))
=-\E_{q(L_i^A)}\operatorname{KL}
\{q(L_i^R)\Vert p_\theta(L_i^R\mid L_i^A,X_i)\}.
\label{eq:missinggap}
\end{equation}
This may penalize dependence solely because the variational family cannot
represent it. It is a valid lower-bound gap, but there is no observed RNA
evidence that warrants it. Integrating out the branch removes this avoidable
gap. Equivalently, an unrestricted conditional variational factor could set
\(q(L_i^R\mid L_i^A)=p_\theta(L_i^R\mid L_i^A,X_i)\).

For each row, retain nodes with observed descendants and eliminate the
remaining terminal branches in reverse generative order. In the implemented
ATAC-to-RNA graph, RNA-only cells retain their ATAC parents; ATAC-only cells
omit all RNA nodes from the fitting objective. Regularization for an omitted
node is omitted for that row as well. Posterior prediction of eliminated
nodes follows the generative ordering after the fit.

\subsection{The same learning equations, two observation models}
For node \(L^m_{ik}\), compute \(a_{ik},b_{ik}\) from modality \(m\)'s mask,
precision, residual and feature moments. Then apply \eqref{eq:learning}
with its active within- and between-modality children. Update \(F^m_k\)
through the original ash normal-means solver using that modality's observed
entries. Missing-modality predictions are never inserted as observations
in a later feature or noise fit.

Removing cross-modality edges gives separate conditional-loading models.
Removing \emph{all} latent-covariate edges also gives the original scalar
cEBMF normal-means updates. Preserving the original code path when no graph
is required retains its initialization, options and finite optimizer behavior.

\subsection{Identification and evaluation}
Paired cells anchor the relationship between modalities. With no pairing,
an unrestricted conditional distribution is generally not identified by the
two marginal distributions; additional assumptions or anchors are needed.
The paired subset should cover the states in which predictions will be used.
Its mere existence is a software precondition, not a statistical sufficiency
guarantee.

Evaluate denoising on paired cells separately from prediction of an entire
missing modality. Hold out complete cell-modality rows to assess the latter.
Fit preprocessing and choose hyperparameters without those held-out values.
The Gaussian derivation applies to the toy observations and appropriately
processed real measurements; a raw-count likelihood requires a separate
extension of the normal-means calculation.
